\documentclass[11pt,a4paper]{article}
\usepackage[utf8]{inputenc}
\usepackage[T1]{fontenc}
\usepackage[catalan]{babel}
\usepackage{lmodern}
\usepackage[sfdefault,lining]{FiraSans}
\usepackage[a4paper,top=2.3cm,bottom=2.2cm,left=1.8cm,right=1.8cm,headheight=1.3cm,headsep=0.4cm,footskip=1cm]{geometry}
\usepackage{amsmath,amssymb}
\usepackage{icomma}
\usepackage{graphicx}
\usepackage[table]{xcolor}
\usepackage{tikz}
\usepackage{fontawesome5}
\usepackage{ccicons}
\usepackage{fancyhdr}
\usepackage{lastpage}
\usepackage{tasks}
\usepackage{needspace}
\usepackage[most]{tcolorbox}
\usepackage{array}
\usepackage[hidelinks,pdftitle={Fitxa d'exercicis de la unitat 9: Probabilitat. 2n de Batxillerat CCSS},pdfsubject={Matemàtiques Aplicades a les Ciències Socials II},pdfkeywords={probabilitat, espai mostral, successos, regla de Laplace, probabilitat condicionada, independència, taules de contingència, diagrames d'arbre, extraccions, probabilitat total, teorema de Bayes, PAU}]{hyperref}

\definecolor{navy}{HTML}{1B2748}
\definecolor{accent}{HTML}{E8603A}
\definecolor{teal}{HTML}{1F9E90}
\definecolor{soft}{HTML}{EEF1F8}
\definecolor{muted}{HTML}{5D6685}

\setlength{\parindent}{0pt}
\setlength{\parskip}{3pt}
\settasks{label=\alph*), label-format=\bfseries\color{navy}, label-width=1.3em, item-indent=1.8em, column-sep=0.6em, after-item-skip=5pt, before-skip=3pt, after-skip=4pt}

% ---------- capçalera i peu ----------
\newcommand{\llicencia}{\href{https://creativecommons.org/licenses/by-nc-sa/4.0/deed.ca}{\ccbyncsa}\ \ Llicència Creative Commons Reconeixement-NoComercial-CompartirIgual 4.0 Internacional (CC BY-NC-SA 4.0)}
\pagestyle{fancy}
\fancyhf{}
\fancyhead[L]{\raisebox{-0.35\height}{\includegraphics[height=0.95cm]{logo-gva.png}}}
\fancyhead[C]{\small\color{muted}Fitxa d'exercicis · Unitat 9: Probabilitat · 2n Batxillerat CCSS}
\fancyhead[R]{\raisebox{-0.35\height}{\includegraphics[height=0.95cm]{logo-vermellar.png}}}
\fancyfoot[L]{\scriptsize\color{muted}\llicencia}
\fancyfoot[R]{\small\color{muted}\thepage}
\renewcommand{\headrulewidth}{0.4pt}
\renewcommand{\headrule}{\color{muted!40}\hrule height \headrulewidth}
\fancypagestyle{primera}{\fancyhf{}\renewcommand{\headrulewidth}{0pt}
  \fancyfoot[L]{\scriptsize\color{muted}\llicencia}
  \fancyfoot[R]{\small\color{muted}\thepage}}

% ---------- nivells i exercicis ----------
\newcommand{\estrelles}[1]{{\color{accent}\ifcase#1\or\faStar\faStar[regular]\faStar[regular]\or\faStar\faStar\faStar[regular]\or\faStar\faStar\faStar\fi}}
\newcounter{exer}
\newcommand{\ex}[1]{\par\needspace{4\baselineskip}\medskip\refstepcounter{exer}\noindent{\color{navy}\bfseries\theexer.}\ \ }

\newcommand{\seccio}[2]{\par\needspace{16\baselineskip}\bigskip
  \begin{tcolorbox}[colback=navy,colframe=navy,coltext=white,arc=2mm,boxrule=0pt,left=3mm,right=3mm,top=1.2mm,bottom=1.2mm]
  {\large\bfseries #1\quad #2}
  \end{tcolorbox}\vspace{-2pt}}

\newcommand{\exf}[1]{\par\medskip\refstepcounter{exer}\noindent{\color{navy}\bfseries\theexer.}\ \ }

% quadre de recordatori (només a la fitxa d'exercicis)
\newcommand{\recorda}[1]{\par\smallskip\begin{tcolorbox}[colback=accent!7,colframe=accent!55,arc=2mm,boxrule=0.5pt,left=2.5mm,right=2.5mm,top=1mm,bottom=1mm]\small{\color{accent}\faLightbulb[regular]}\ \textbf{Recorda:}\ #1\end{tcolorbox}}
% etiqueta dels exercicis de la PAU d'altres comunitats
\newcommand{\pau}[1]{{\color{teal}\small\bfseries\faUniversity\ PAU · #1.}\ }
% successos contraris i probabilitat condicionada
\newcommand{\co}[1]{\overline{#1}}
\newcommand{\cond}{\,|\,}
\newcommand{\res}[1]{#1}\newcommand{\rest}[1]{#1}

\def\fmneed{11}% que el títol d'una secció no quede sol al final del full
% =====================================================================
% Capçalera FentMates de les fitxes: un color per a cada curs.
% Els colors són els de la pàgina inicial (LEVELS a assets/js/materials.js);
% el logo ve de _eines/logo/pdf/ (python3 _eines/logo/genera.py).
%
% Ús, al final del preàmbul de la unitat:
%   % =====================================================================
% Capçalera FentMates de les fitxes: un color per a cada curs.
% Els colors són els de la pàgina inicial (LEVELS a assets/js/materials.js);
% el logo ve de _eines/logo/pdf/ (python3 _eines/logo/genera.py).
%
% Ús, al final del preàmbul de la unitat:
%   % =====================================================================
% Capçalera FentMates de les fitxes: un color per a cada curs.
% Els colors són els de la pàgina inicial (LEVELS a assets/js/materials.js);
% el logo ve de _eines/logo/pdf/ (python3 _eines/logo/genera.py).
%
% Ús, al final del preàmbul de la unitat:
%   \input{../comu/fentmates.tex}
%   \fmcurs{1eso}
% i, al començament del document:
%   \capcalera{Matemàtiques}{1r ESO}{Fitxa d'exercicis · Unitat 2}{Divisibilitat}
% Opcional abans de \input: \def\fmneed{14} (espai mínim abans de cada secció).
% =====================================================================
\makeatletter
\def\fm@logodir{../../logo/pdf}
\providecommand{\fmneed}{10}
\def\fm@cursnom{}

% clau del curs, color i color del text damunt del color
\newcommand{\fm@curs}[3]{\@namedef{fm@color@#1}{#2}\@namedef{fm@text@#1}{#3}}
\fm@curs{1eso}{E6B84A}{navy}
\fm@curs{2eso}{3B82F6}{white}
\fm@curs{3eso}{FF7A50}{white}
\fm@curs{4eso-a}{A78BFA}{white}
\fm@curs{4eso-b}{6D28D9}{white}
\fm@curs{1batx-mi}{9B1C1C}{white}
\fm@curs{1batx-ccss}{EF6B6B}{white}
\fm@curs{2batx-mii}{1E6B47}{white}
\fm@curs{2batx-ccss}{2BB3A3}{white}

\newcommand{\fmcurs}[1]{%
  \@ifundefined{fm@color@#1}{\PackageError{fentmates}{Curs desconegut: #1}{}}{}%
  \edef\fm@tmp{\noexpand\definecolor{curs}{HTML}{\@nameuse{fm@color@#1}}%
               \noexpand\colorlet{curstext}{\@nameuse{fm@text@#1}}}\fm@tmp
  \colorlet{cursfosc}{curs!72!black}%
  \def\fm@logo{\fm@logodir/vertical-#1.pdf}%
  \def\fm@marca{\fm@logodir/marca-#1.pdf}}

% ---------- capçalera de la primera pàgina ----------
% #1 matèria (petit), #2 curs, #3 tipus de document i unitat, #4 títol
\newcommand{\capcalera}[4]{%
  \gdef\fm@cursnom{#2}%
  \par\noindent
  \begin{tcolorbox}[colback=curs,colframe=curs,boxrule=0pt,arc=2.2mm,width=3.5cm,height=2.15cm,
      valign=center,left=3mm,right=2mm,top=1mm,bottom=1mm,nobeforeafter,box align=center]
    \color{curstext}{\footnotesize #1}\par\vspace{1pt}{\fontsize{19}{21}\selectfont\bfseries #2\par}%
  \end{tcolorbox}\hspace{5mm}%
  \begin{minipage}[c]{\dimexpr\linewidth-3.5cm-5mm-2.9cm\relax}
    {\color{cursfosc}\small\bfseries\MakeUppercase{#3}}\par\vspace{2pt}
    {\color{navy}\fontsize{20}{23}\selectfont\bfseries #4\par}
  \end{minipage}\hfill
  \begin{minipage}[c]{2.6cm}\raggedleft\includegraphics[width=2.6cm]{\fm@logo}\end{minipage}\par
  \vspace{2pt}}

% ---------- barres de secció del color del curs ----------
\renewcommand{\seccio}[2]{\par\needspace{\fmneed\baselineskip}\bigskip
  \begin{tcolorbox}[colback=curs,colframe=curs,coltext=curstext,arc=2mm,boxrule=0pt,left=3mm,right=3mm,top=1.2mm,bottom=1.2mm]
  {\large\bfseries #1\quad #2}
  \end{tcolorbox}\vspace{-2pt}}

% ---------- capçalera i peu de les pàgines ----------
% primera pàgina: logos institucionals al peu
\fancypagestyle{primera}{\fancyhf{}\renewcommand{\headrulewidth}{0pt}
  \fancyfoot[L]{\raisebox{-0.32\height}{\includegraphics[height=0.78cm]{logo-gva.png}}\hspace{5mm}%
    \raisebox{-0.32\height}{\includegraphics[height=0.78cm]{logo-vermellar.png}}}
  \fancyfoot[R]{\scriptsize\color{muted}\href{https://creativecommons.org/licenses/by-nc-sa/4.0/deed.ca}{\ccbyncsa}\ \ CC BY-NC-SA 4.0 · fentmates.org\qquad\small\thepage}}
% la resta: el símbol a l'esquerra i el curs a la dreta, amb el títol al centre com abans
\fancyhead[L]{\raisebox{-0.3\height}{\includegraphics[height=0.62cm]{\fm@marca}}}
\fancyhead[R]{\small\bfseries\color{cursfosc}\fm@cursnom}
\renewcommand{\headrule}{\color{curs!45}\hrule height \headrulewidth}
\makeatother

%   \fmcurs{1eso}
% i, al començament del document:
%   \capcalera{Matemàtiques}{1r ESO}{Fitxa d'exercicis · Unitat 2}{Divisibilitat}
% Opcional abans de \input: \def\fmneed{14} (espai mínim abans de cada secció).
% =====================================================================
\makeatletter
\def\fm@logodir{../../logo/pdf}
\providecommand{\fmneed}{10}
\def\fm@cursnom{}

% clau del curs, color i color del text damunt del color
\newcommand{\fm@curs}[3]{\@namedef{fm@color@#1}{#2}\@namedef{fm@text@#1}{#3}}
\fm@curs{1eso}{E6B84A}{navy}
\fm@curs{2eso}{3B82F6}{white}
\fm@curs{3eso}{FF7A50}{white}
\fm@curs{4eso-a}{A78BFA}{white}
\fm@curs{4eso-b}{6D28D9}{white}
\fm@curs{1batx-mi}{9B1C1C}{white}
\fm@curs{1batx-ccss}{EF6B6B}{white}
\fm@curs{2batx-mii}{1E6B47}{white}
\fm@curs{2batx-ccss}{2BB3A3}{white}

\newcommand{\fmcurs}[1]{%
  \@ifundefined{fm@color@#1}{\PackageError{fentmates}{Curs desconegut: #1}{}}{}%
  \edef\fm@tmp{\noexpand\definecolor{curs}{HTML}{\@nameuse{fm@color@#1}}%
               \noexpand\colorlet{curstext}{\@nameuse{fm@text@#1}}}\fm@tmp
  \colorlet{cursfosc}{curs!72!black}%
  \def\fm@logo{\fm@logodir/vertical-#1.pdf}%
  \def\fm@marca{\fm@logodir/marca-#1.pdf}}

% ---------- capçalera de la primera pàgina ----------
% #1 matèria (petit), #2 curs, #3 tipus de document i unitat, #4 títol
\newcommand{\capcalera}[4]{%
  \gdef\fm@cursnom{#2}%
  \par\noindent
  \begin{tcolorbox}[colback=curs,colframe=curs,boxrule=0pt,arc=2.2mm,width=3.5cm,height=2.15cm,
      valign=center,left=3mm,right=2mm,top=1mm,bottom=1mm,nobeforeafter,box align=center]
    \color{curstext}{\footnotesize #1}\par\vspace{1pt}{\fontsize{19}{21}\selectfont\bfseries #2\par}%
  \end{tcolorbox}\hspace{5mm}%
  \begin{minipage}[c]{\dimexpr\linewidth-3.5cm-5mm-2.9cm\relax}
    {\color{cursfosc}\small\bfseries\MakeUppercase{#3}}\par\vspace{2pt}
    {\color{navy}\fontsize{20}{23}\selectfont\bfseries #4\par}
  \end{minipage}\hfill
  \begin{minipage}[c]{2.6cm}\raggedleft\includegraphics[width=2.6cm]{\fm@logo}\end{minipage}\par
  \vspace{2pt}}

% ---------- barres de secció del color del curs ----------
\renewcommand{\seccio}[2]{\par\needspace{\fmneed\baselineskip}\bigskip
  \begin{tcolorbox}[colback=curs,colframe=curs,coltext=curstext,arc=2mm,boxrule=0pt,left=3mm,right=3mm,top=1.2mm,bottom=1.2mm]
  {\large\bfseries #1\quad #2}
  \end{tcolorbox}\vspace{-2pt}}

% ---------- capçalera i peu de les pàgines ----------
% primera pàgina: logos institucionals al peu
\fancypagestyle{primera}{\fancyhf{}\renewcommand{\headrulewidth}{0pt}
  \fancyfoot[L]{\raisebox{-0.32\height}{\includegraphics[height=0.78cm]{logo-gva.png}}\hspace{5mm}%
    \raisebox{-0.32\height}{\includegraphics[height=0.78cm]{logo-vermellar.png}}}
  \fancyfoot[R]{\scriptsize\color{muted}\href{https://creativecommons.org/licenses/by-nc-sa/4.0/deed.ca}{\ccbyncsa}\ \ CC BY-NC-SA 4.0 · fentmates.org\qquad\small\thepage}}
% la resta: el símbol a l'esquerra i el curs a la dreta, amb el títol al centre com abans
\fancyhead[L]{\raisebox{-0.3\height}{\includegraphics[height=0.62cm]{\fm@marca}}}
\fancyhead[R]{\small\bfseries\color{cursfosc}\fm@cursnom}
\renewcommand{\headrule}{\color{curs!45}\hrule height \headrulewidth}
\makeatother

%   \fmcurs{1eso}
% i, al començament del document:
%   \capcalera{Matemàtiques}{1r ESO}{Fitxa d'exercicis · Unitat 2}{Divisibilitat}
% Opcional abans de \input: \def\fmneed{14} (espai mínim abans de cada secció).
% =====================================================================
\makeatletter
\def\fm@logodir{../../logo/pdf}
\providecommand{\fmneed}{10}
\def\fm@cursnom{}

% clau del curs, color i color del text damunt del color
\newcommand{\fm@curs}[3]{\@namedef{fm@color@#1}{#2}\@namedef{fm@text@#1}{#3}}
\fm@curs{1eso}{E6B84A}{navy}
\fm@curs{2eso}{3B82F6}{white}
\fm@curs{3eso}{FF7A50}{white}
\fm@curs{4eso-a}{A78BFA}{white}
\fm@curs{4eso-b}{6D28D9}{white}
\fm@curs{1batx-mi}{9B1C1C}{white}
\fm@curs{1batx-ccss}{EF6B6B}{white}
\fm@curs{2batx-mii}{1E6B47}{white}
\fm@curs{2batx-ccss}{2BB3A3}{white}

\newcommand{\fmcurs}[1]{%
  \@ifundefined{fm@color@#1}{\PackageError{fentmates}{Curs desconegut: #1}{}}{}%
  \edef\fm@tmp{\noexpand\definecolor{curs}{HTML}{\@nameuse{fm@color@#1}}%
               \noexpand\colorlet{curstext}{\@nameuse{fm@text@#1}}}\fm@tmp
  \colorlet{cursfosc}{curs!72!black}%
  \def\fm@logo{\fm@logodir/vertical-#1.pdf}%
  \def\fm@marca{\fm@logodir/marca-#1.pdf}}

% ---------- capçalera de la primera pàgina ----------
% #1 matèria (petit), #2 curs, #3 tipus de document i unitat, #4 títol
\newcommand{\capcalera}[4]{%
  \gdef\fm@cursnom{#2}%
  \par\noindent
  \begin{tcolorbox}[colback=curs,colframe=curs,boxrule=0pt,arc=2.2mm,width=3.5cm,height=2.15cm,
      valign=center,left=3mm,right=2mm,top=1mm,bottom=1mm,nobeforeafter,box align=center]
    \color{curstext}{\footnotesize #1}\par\vspace{1pt}{\fontsize{19}{21}\selectfont\bfseries #2\par}%
  \end{tcolorbox}\hspace{5mm}%
  \begin{minipage}[c]{\dimexpr\linewidth-3.5cm-5mm-2.9cm\relax}
    {\color{cursfosc}\small\bfseries\MakeUppercase{#3}}\par\vspace{2pt}
    {\color{navy}\fontsize{20}{23}\selectfont\bfseries #4\par}
  \end{minipage}\hfill
  \begin{minipage}[c]{2.6cm}\raggedleft\includegraphics[width=2.6cm]{\fm@logo}\end{minipage}\par
  \vspace{2pt}}

% ---------- barres de secció del color del curs ----------
\renewcommand{\seccio}[2]{\par\needspace{\fmneed\baselineskip}\bigskip
  \begin{tcolorbox}[colback=curs,colframe=curs,coltext=curstext,arc=2mm,boxrule=0pt,left=3mm,right=3mm,top=1.2mm,bottom=1.2mm]
  {\large\bfseries #1\quad #2}
  \end{tcolorbox}\vspace{-2pt}}

% ---------- capçalera i peu de les pàgines ----------
% primera pàgina: logos institucionals al peu
\fancypagestyle{primera}{\fancyhf{}\renewcommand{\headrulewidth}{0pt}
  \fancyfoot[L]{\raisebox{-0.32\height}{\includegraphics[height=0.78cm]{logo-gva.png}}\hspace{5mm}%
    \raisebox{-0.32\height}{\includegraphics[height=0.78cm]{logo-vermellar.png}}}
  \fancyfoot[R]{\scriptsize\color{muted}\href{https://creativecommons.org/licenses/by-nc-sa/4.0/deed.ca}{\ccbyncsa}\ \ CC BY-NC-SA 4.0 · fentmates.org\qquad\small\thepage}}
% la resta: el símbol a l'esquerra i el curs a la dreta, amb el títol al centre com abans
\fancyhead[L]{\raisebox{-0.3\height}{\includegraphics[height=0.62cm]{\fm@marca}}}
\fancyhead[R]{\small\bfseries\color{cursfosc}\fm@cursnom}
\renewcommand{\headrule}{\color{curs!45}\hrule height \headrulewidth}
\makeatother

\fmcurs{2batx-ccss}

\begin{document}
\thispagestyle{primera}

\capcalera{Matemàtiques CCSS}{2n Batx}{Fitxa d'exercicis · Unitat 9}{\raggedright Probabilitat}
\vspace{10pt}
Nom i cognoms: \rule{8.4cm}{0.4pt}\hfill Grup: \rule{1.6cm}{0.4pt}\hfill Data: \rule{2.2cm}{0.4pt}
\vspace{8pt}

\begin{tcolorbox}[colback=soft,colframe=soft,arc=2mm,left=3mm,right=3mm,top=1.5mm,bottom=1.5mm]
\small Resol els exercicis al quadern mostrant tots els passos. Defineix sempre els successos amb lletres i escriu les dades amb la notació de probabilitats ($P(A)$, $P(A\cap B)$, $P(A\,|\,B)$\dots). Fes un diagrama de Venn, una taula de contingència o un diagrama d'arbre quan t'ajude. Els 50 exercicis marcats amb \textcolor{teal}{\faUniversity\ PAU} són de les proves d'accés a la universitat d'altres comunitats autònomes.
\end{tcolorbox}

% =====================================================
\seccio{1}{Espai mostral i operacions amb successos}

\exf{1} Escriu l'espai mostral dels experiments següents i digues quants elements té:\begin{tasks}(1)
  \task Llançar una moneda i un dau.
  \task Anotar, per ordre de naixement, si cadascun dels tres fills d'una família és xic ($X$) o xica ($Y$).
  \task Traure, una darrere l'altra i sense reemplaçament, dues boles d'una urna que en té una de roja ($R$), una de blava ($B$) i una de verda ($V$).
  \task Llançar un dau fins que isca un 6 i anotar el nombre de tirades que han calgut.
\end{tasks}

\ex{1} Llancem un dau i considerem els successos $A=$«eixir un nombre menor que 4», $B=$«eixir parell» i $C=$«eixir 5 o 6». Calcula:\begin{tasks}(4)
  \task $A\cup B$
  \task $A\cap B$
  \task $\co A$
  \task $\co B$
  \task $A-B$
  \task $B\cap C$
  \task $\co A\cap\co B$
  \task $\overline{A\cup C}$
\end{tasks}Quins parells d'aquests successos són incompatibles?

\ex{1} Siguen $M=$«un alumne aprova Matemàtiques» i $N=$«aprova Anglés». Expressa amb operacions entre $M$ i $N$ els successos següents:\begin{tasks}(2)
  \task Aprova les dues.
  \task No n'aprova cap.
  \task Aprova només Matemàtiques.
  \task Aprova com a molt una.
  \task Aprova almenys una.
  \task Aprova exactament una.
  \task Suspén alguna de les dues.
\end{tasks}\recorda{«almenys un» és la unió; «tots dos», la intersecció; «cap» és $\co M\cap\co N=\overline{M\cup N}$ (De Morgan).}

\ex{1} En l'espai mostral $E=\{1,2,3,4,5,6,7,8\}$ considerem $A=\{1,2,3,4\}$ i $B=\{2,4,6,8\}$. Comprova les dues lleis de De Morgan.

\ex{2} Siguen $A$, $B$ i $C$ tres successos d'un mateix experiment. Expressa amb operacions:\begin{tasks}(2)
  \task Ocorre només $A$.
  \task Ocorre almenys un dels tres.
  \task No n'ocorre cap.
  \task N'ocorre exactament un.
  \task Ocorren tots tres.
  \task N'ocorren almenys dos.
\end{tasks}


% =====================================================
\seccio{2}{Probabilitat. Regla de Laplace}

\exf{1} Una urna conté 4 boles roges, 5 de blaves i 3 de verdes. Traiem una bola a l'atzar. Calcula la probabilitat que:\begin{tasks}(4)
  \task siga roja;
  \task no siga verda;
  \task siga roja o blava;
  \task siga groga.
\end{tasks}

\ex{1} Traiem una carta d'una baralla espanyola de 40 cartes (quatre pals de deu cartes: de l'1 al 7, sota, cavall i rei). Calcula la probabilitat de traure:\begin{tasks}(2)
  \task un or;
  \task una figura (sota, cavall o rei);
  \task un or o una figura;
  \task una carta que no siga ni or ni figura;
  \task un as que no siga d'ors.
\end{tasks}

\ex{1} \pau{Castella i Lleó, EBAU ordinària 2020} Es llança una moneda tres vegades. Calcula la probabilitat d'obtindre almenys una creu.

\ex{1} \pau{Castella i Lleó, EBAU extraordinària 2024} Una botiga d'animals fa un sorteig amb paperetes de tres xifres. El número premiat s'obté traient a l'atzar cada xifra, per separat i amb reemplaçament, d'una bossa amb boles del 0 al 9. Calcula la probabilitat que el número premiat acabe en 55.

\ex{1} Llancem dos daus. Calcula la probabilitat que:\begin{tasks}(2)
  \task la suma siga 9;
  \task la suma siga com a molt 4;
  \task el producte siga 12;
  \task els dos resultats siguen parells;
  \task la diferència entre els dos resultats siga almenys 3.
\end{tasks}

\ex{2} \pau{Illes Balears, PBAU juliol 2018} Un dau està carregat de manera que la probabilitat d'obtindre un 6 és $\frac12$ i les probabilitats d'obtindre cadascuna de les altres cares són iguals a $p$. Es llança aquest dau; calcula la probabilitat de cada un dels successos següents:\begin{tasks}(2)
  \task S'obté un dos.
  \task No s'obté un tres.
  \task S'obté un nombre parell.
  \task S'obté un nombre senar.
\end{tasks}

\ex{2} \pau{Andalusia, PEvAU ordinària 2023} Tenim una moneda trucada en què la probabilitat d'obtindre cara és el doble que la d'obtindre creu.\begin{tasks}(1)
  \task Calcula la probabilitat d'obtindre cara en llançar la moneda.
  \task Calcula la probabilitat d'obtindre una cara i una creu, sense importar l'ordre, en llançar dues vegades la moneda.
  \task Calcula la probabilitat d'obtindre almenys una cara en llançar dues vegades la moneda.
\end{tasks}

\ex{1} Triem a l'atzar un nombre de l'1 al 100. Calcula la probabilitat que siga múltiple de 3 o de 5.

\ex{1} Llancem tres monedes. Calcula la probabilitat d'obtindre exactament dues cares i la d'obtindre més cares que creus.

\ex{2} Raona si les afirmacions següents són vertaderes o falses:\begin{tasks}(1)
  \task Si llancem dos daus i sumem, la probabilitat de cada suma és $\frac1{11}$, perquè hi ha 11 sumes possibles.
  \task Si una moneda equilibrada ha eixit cara cinc vegades seguides, és més probable que la sisena vegada isca creu.
  \task Pot ser que $P(A)=1{,}2$ si $A$ és un succés molt probable.
\end{tasks}


% =====================================================
\seccio{3}{Fórmules de la probabilitat}

\exf{1} Siguen $A$ i $B$ dos successos amb $P(A)=0{,}45$, $P(B)=0{,}3$ i $P(A\cup B)=0{,}6$. Calcula:\begin{tasks}(3)
  \task $P(A\cap B)$
  \task $P(\co A\cap\co B)$
  \task $P(A\cap\co B)$
  \task $P(\co A\cup\co B)$
  \task $P(\text{només un dels dos})$
\end{tasks}\recorda{$P(A\cup B)=P(A)+P(B)-P(A\cap B)$; \ $P(\co A)=1-P(A)$; \ $P(A\cap\co B)=P(A)-P(A\cap B)$; \ De Morgan: $\co A\cap\co B=\overline{A\cup B}$ i $\co A\cup\co B=\overline{A\cap B}$.}

\ex{1} \pau{Madrid, EvAU juliol 2018} Siguen $A$ i $B$ dos successos d'un experiment aleatori tals que $P(A)=0{,}4$, $P(B)=0{,}6$ i $P(A\cup B)=0{,}8$. Calcula $P(\co A\cap B)$ i $P(\co A\cup\co B)$.

\ex{1} De dos successos se sap que $P(\co A)=0{,}7$, $P(\co B)=0{,}4$ i $P(\co A\cup\co B)=0{,}9$. Calcula $P(A)$, $P(B)$, $P(A\cap B)$, $P(A\cup B)$ i $P(\co A\cap\co B)$.

\ex{1} Els successos $A$ i $B$ són incompatibles, amb $P(A)=0{,}35$ i $P(B)=0{,}4$. Calcula $P(A\cup B)$, $P(\co A\cap\co B)$, $P(A\cap\co B)$ i $P(\co A\cup\co B)$.

\ex{2} Siguen $A$ i $B$ dos successos amb $P(A)=0{,}6$ i $P(B)=0{,}7$.\begin{tasks}(1)
  \task Poden ser incompatibles?
  \task Entre quins valors pot estar $P(A\cap B)$?
  \task I $P(A\cup B)$?
\end{tasks}

\ex{2} Tres successos $A$, $B$ i $C$ tenen probabilitat $0{,}4$ cadascun; les interseccions de dos en dos tenen probabilitat $0{,}15$ i $P(A\cap B\cap C)=0{,}05$. Calcula la probabilitat que n'ocórrega almenys un i que no n'ocórrega cap.


% =====================================================
\seccio{4}{Càlcul de probabilitats coneixent l'espai mostral i sense conéixer-lo}

\exf{1} Llancem dos daus i considerem els successos $A=$«la suma és parella» i $B=$«almenys un dels daus és un 6». Calcula $P(A)$, $P(B)$ i $P(A\cap B)$ per recompte i, amb la fórmula, $P(A\cup B)$. Comprova el resultat comptant els casos de $A\cup B$.

\ex{1} Traiem a l'atzar una targeta d'una caixa amb targetes numerades de l'1 al 30. Siguen $A=$«el nombre és parell» i $B=$«el nombre és múltiple de 3». Calcula, aplicant les fórmules, la probabilitat que siga parell o múltiple de 3, que complisca només una de les dues condicions i que no en complisca cap.

\ex{1} \pau{Galícia, ABAU extraordinària 2024} En una enquesta, el 80\,\% dels entrevistats diu que llig o escolta música, el 35\,\% fa les dues coses i el 60\,\% no llig. Calcula les probabilitats que una persona triada a l'atzar:\begin{tasks}(1)
  \task escolte música i no llija;
  \task llija i no escolte música;
  \task faça només una de les dues coses.
\end{tasks}

\ex{1} \pau{Andalusia, PEvAU ordinària 2022} En el seu temps lliure, el 65\,\% dels estudiants d'un centre juga a videojocs, el 45\,\% llig llibres i el 15\,\% no fa cap de les dues coses. Triem un estudiant a l'atzar. Calcula la probabilitat que:\begin{tasks}(1)
  \task jugue a videojocs o llija llibres;
  \task jugue a videojocs i no llija llibres.
\end{tasks}

\ex{1} En una ciutat, el 55\,\% dels habitants llig la premsa, el 40\,\% veu les notícies a la televisió i el 25\,\% no fa cap de les dues coses. Construeix la taula de contingència i calcula la probabilitat que un habitant triat a l'atzar:\begin{tasks}(2)
  \task faça les dues coses;
  \task només llija la premsa;
  \task faça només una de les dues;
  \task llija la premsa o no veja la televisió.
\end{tasks}

\ex{1} En un grup de 30 persones, 18 parlen anglés, 12 francés i 6 no parlen cap dels dos idiomes. Triem una persona a l'atzar. Calcula la probabilitat que parle els dos idiomes, que en parle només un i que parle francés però no anglés.


% =====================================================
\seccio{5}{Probabilitat condicionada}

\exf{1} Siguen $A$ i $B$ dos successos amb $P(A)=0{,}5$, $P(B)=0{,}4$ i $P(A\cap B)=0{,}1$. Calcula:\begin{tasks}(3)
  \task $P(A\cond B)$
  \task $P(B\cond A)$
  \task $P(A\cond\co B)$
  \task $P(\co A\cond\co B)$
  \task $P(B\cond\co A)$
  \task $P(\co B\cond A)$
\end{tasks}\recorda{$P(A\cond B)=\dfrac{P(A\cap B)}{P(B)}$. El succés que va darrere de la barra és la condició: va al denominador.}

\ex{2} \pau{Madrid, EvAU juliol 2019} Siguen $A$ i $B$ dos successos amb $P(A)=0{,}3$, $P(B\cond A)=0{,}4$ i $P(B\cond\co A)=0{,}6$. Calcula:\begin{tasks}(2)
  \task $P(A\cond B)$
  \task $P(\co A\cond\co B)$
\end{tasks}

\ex{1} \pau{Illes Balears, PBAU setembre 2017} Siguen $A$ i $B$ dos successos amb probabilitats $0{,}4$ i $0{,}6$, respectivament. Se sap que, donat $B$, la probabilitat que ocórrega $A$ és $0{,}3$.\begin{tasks}(1)
  \task Quina és la probabilitat que ocórreguen tots dos successos alhora?
  \task Quina és la probabilitat que ocórrega algun dels dos?
  \task Quina és la probabilitat que no n'ocórrega cap?
\end{tasks}

\ex{1} \pau{Andalusia, PEvAU juny 2017} El 90\,\% dels alumnes d'un centre està interessat per les xarxes socials, el 60\,\% per les seues notes i el 55\,\% per les dues coses. Es tria un alumne a l'atzar.\begin{tasks}(1)
  \task Quina és la probabilitat que estiga interessat per alguna de les dues coses?
  \task Calcula la probabilitat que estiga interessat per les notes, sabent que no està interessat per les xarxes socials.
  \task Calcula la probabilitat que no estiga interessat per cap de les dues coses.
\end{tasks}

\ex{2} \pau{Andalusia, PEvAU ordinària 2020} A 120 estudiants se'ls ha recomanat la lectura de dos llibres: 46 han llegit el primer, 34 el segon i 16 tots dos. Es tria un estudiant a l'atzar. Calcula la probabilitat que:\begin{tasks}(1)
  \task haja llegit algun dels dos llibres;
  \task no n'haja llegit cap;
  \task només haja llegit el primer;
  \task haja llegit el primer, si se sap que no ha llegit el segon.
\end{tasks}

\ex{2} \pau{Andalusia, PEvAU ordinària 2021} De les compres fetes en les últimes rebaixes, el 55\,\% van ser de productes electrònics i el 72\,\% es van fer per Internet; de les compres fetes per Internet, el 64\,\% van ser de productes electrònics. Es tria una compra a l'atzar. Calcula la probabilitat que:\begin{tasks}(1)
  \task haja sigut de productes electrònics i feta per Internet;
  \task s'haja fet per Internet o haja sigut de productes electrònics;
  \task no s'haja fet per Internet, sabent que no va ser de productes electrònics.
\end{tasks}

\ex{2} \pau{Múrcia, PAU extraordinària 2025} En una enquesta a 100 estudiants de segon de Batxillerat, 30 volen estudiar Veterinària, 25 Economia, 20 Química, 15 Biologia i 10 Matemàtiques. D'ells, 10, 12, 13, 8 i 6 són dones, respectivament. Ningú no es pot matricular en més d'un grau. Triem un estudiant a l'atzar. Calcula la probabilitat que:\begin{tasks}(2)
  \task siga dona;
  \task siga home i vulga estudiar Veterinària;
  \task vulga estudiar Química o Veterinària;
  \task no vulga estudiar ni Biologia ni Matemàtiques;
  \task vulga estudiar Química, sabent que és dona;
  \task vulga estudiar Economia, sabent que és home.
\end{tasks}

\ex{2} De dos successos se sap que $P(A)=0{,}6$, $P(B\cond A)=0{,}5$ i $P(B\cond\co A)=0{,}25$. Calcula $P(A\cap B)$, $P(B)$, $P(A\cond B)$ i $P(A\cup B)$.


% =====================================================
\seccio{6}{Independència de successos}

\exf{1} Els successos $A$ i $B$ són independents, amb $P(A)=0{,}3$ i $P(B)=0{,}5$. Calcula $P(A\cap B)$, $P(A\cup B)$, $P(\co A\cap\co B)$, $P(A\cond B)$ i $P(A\cap\co B)$.\recorda{$A$ i $B$ són independents si $P(A\cap B)=P(A)\cdot P(B)$. Aleshores també ho són $A$ i $\co B$, $\co A$ i $B$, i $\co A$ i $\co B$.}

\ex{2} Se sap que $P(A)=0{,}4$ i $P(A\cup B)=0{,}7$. Calcula $P(B)$ si:\begin{tasks}(2)
  \task $A$ i $B$ són incompatibles;
  \task $A$ i $B$ són independents.
\end{tasks}

\ex{1} \pau{Madrid, EvAU extraordinària 2021} Siguen $A$ i $B$ dos successos amb $P(A)=0{,}5$, $P(B)=0{,}8$ i $P(A\cup B)=0{,}9$.\begin{tasks}(1)
  \task Estudia si els successos $A$ i $B$ són independents.
  \task Calcula $P(A\cond\co B)$.
\end{tasks}

\ex{1} \pau{Castella-la Manxa, EvAU setembre 2017} D'un estudi sobre accidents de trànsit es va deduir que el 29\,\% dels conductors superaven els límits d'alcohol en sang, el 14\,\% tenien presència de drogues en sang i el 37\,\% superaven els límits d'alcohol o tenien presència de drogues (o totes dues coses).\begin{tasks}(1)
  \task Calcula la probabilitat que, en un accident, el conductor supere els límits d'alcohol i tinga presència de drogues en sang.
  \task Raona si són independents els successos «superar els límits d'alcohol» i «presència de drogues en sang».
\end{tasks}

\ex{2} \pau{Galícia, ABAU ordinària 2020} Siguen $A$ i $B$ dos successos amb $P(A)=0{,}4$, $P(B)=0{,}7$ i $P(B\cond A)=0{,}75$. Calcula:\begin{tasks}(4)
  \task $P(A\cap B)$
  \task $P(A\cup\co B)$
  \task $P(A\cap\co B)$
  \task Són independents?
\end{tasks}

\ex{1} \pau{Andalusia, PEvAU setembre 2018} En una localitat, el 25\,\% dels habitants va periòdicament al dentista, el 10\,\% es fa una analítica i el 8\,\% fa les dues coses.\begin{tasks}(1)
  \task Raona si els successos «anar al dentista» i «fer-se una analítica» són independents.
  \task Quin percentatge d'habitants ni es fa una analítica ni va al dentista?
  \task Si triem a l'atzar un habitant d'entre els que van al dentista, quina és la probabilitat que es faça una analítica?
\end{tasks}

\ex{1} \pau{La Rioja, EBAU ordinària 2020} En una classe hi ha 24 estudiants: 12 han aprovat anglés, 16 han aprovat matemàtiques i 4 han suspés les dues assignatures.\begin{tasks}(1)
  \task Quina és la probabilitat que un alumne triat a l'atzar haja aprovat matemàtiques i suspés anglés?
  \task Quina és la probabilitat que haja aprovat les dues assignatures?
  \task Són independents els successos «aprovar matemàtiques» i «aprovar anglés»?
\end{tasks}

\ex{2} \pau{Galícia, ABAU ordinària 2024} S'estima que en una població el 20\,\% pateix obesitat i que l'11\,\% pateix obesitat i és hipertens. A més, el 27,5\,\% dels hipertensos pateix obesitat.\begin{tasks}(1)
  \task Quin percentatge de la població pateix obesitat o és hipertens?
  \task Són independents els successos «patir obesitat» i «ser hipertens»?
  \task Calcula la probabilitat que un individu que no pateix obesitat siga hipertens.
\end{tasks}

\ex{2} \pau{Galícia, ABAU setembre 2017} El 60\,\% dels individus d'una població estan vacunats contra una malaltia. Durant una epidèmia, el 20\,\% va contraure la malaltia i el 3\,\% està vacunat i la va contraure.\begin{tasks}(1)
  \task Calcula el percentatge d'individus que han contret la malaltia entre els que no estan vacunats.
  \task Calcula el percentatge de vacunats entre els que van contraure la malaltia. Justifica si els successos «estar vacunat» i «contraure la malaltia» són dependents o independents.
\end{tasks}

\ex{2} \pau{Madrid, PAU ordinària 2025} De tres successos $A$, $B$ i $C$ se sap que $A$ i $C$ són disjunts (incompatibles), $A$ i $B$ són independents, $P(A)=0{,}25$, $P(B)=0{,}2$ i $P(B\cap C)=0{,}05$.\begin{tasks}(1)
  \task Calcula la probabilitat que ocórrega almenys un dels successos $A$ o $B$.
  \task Calcula $P(\co B\cup\co C)$.
  \task Poden ser independents els successos $A$ i $C$?
\end{tasks}

\ex{1} \pau{País Basc, EAU extraordinària 2022} La Deiene i la Kattalin són jugadores de bàsquet. La Deiene encistella 2 de cada 5 tirs i la Kattalin 3 de cada 7. Si totes dues tiren a cistella una sola vegada, calcula la probabilitat que:\begin{tasks}(2)
  \task totes dues encistellen;
  \task no n'encistelle cap;
  \task només encistelle la Deiene;
  \task n'encistelle almenys una.
\end{tasks}

\ex{2} \pau{Illes Balears, PBAU extraordinària 2023} Tirem dos daus no trucats i considerem els successos $A=$«en el primer dau ha eixit un 1», $B=$«en el segon dau ha eixit un 1» i $C=$«la suma dels dos daus és 3».\begin{tasks}(1)
  \task Calcula $P(A)$.
  \task Calcula $P(A\cup B)$.
  \task Són independents els successos $C$ i $A\cup B$?
\end{tasks}

\ex{2} Raona si les afirmacions següents són vertaderes o falses:\begin{tasks}(1)
  \task Dos successos incompatibles (amb probabilitats no nul·les) són independents.
  \task Si $A$ i $B$ són independents, aleshores $\co A$ i $\co B$ també ho són.
  \task Si $P(A\cond B)=P(A)$, aleshores $P(B\cond A)=P(B)$.
\end{tasks}


% =====================================================
\seccio{7}{Experiments compostos. Extraccions amb i sense reemplaçament}

\exf{1} Una urna conté 4 boles roges i 6 de blaves. Traiem tres boles. Calcula la probabilitat que les tres siguen roges i la probabilitat que n'hi haja exactament dues de roges:\begin{tasks}(2)
  \task amb reemplaçament;
  \task sense reemplaçament.
\end{tasks}\recorda{en un diagrama d'arbre, multipliquem les probabilitats al llarg de cada camí i sumem els camins que formen el succés. Sense reemplaçament, la composició de l'urna canvia en cada extracció.}

\ex{1} \pau{Aragó, EvAU extraordinària 2020} En una bossa tenim 8 boles: 3 blanques, 1 roja i 4 negres. Traiem dues boles sense reemplaçament. Calcula:\begin{tasks}(1)
  \task La probabilitat que les dues siguen blanques.
  \task La probabilitat que almenys una siga blanca.
  \task La probabilitat que les dues siguen del mateix color.
  \task La probabilitat que siguen blanques, si se sap que les dues són del mateix color.
\end{tasks}

\ex{1} \pau{País Basc, EAU juny 2018} En una urna hi ha 15 boles blanques i 5 de negres. Calcula:\begin{tasks}(1)
  \task Si traiem una bola a l'atzar, la probabilitat que siga blanca.
  \task Si en traiem dues a l'atzar, la probabilitat que totes dues siguen blanques.
  \task Si traiem una bola i després una altra, i la primera és negra, la probabilitat que la segona també siga negra.
  \task Si traiem una bola i després una altra, la probabilitat que siguen de color diferent.
\end{tasks}

\ex{1} \pau{Navarra, EvAU setembre 2017} Una caixa conté dotze bolígrafs, quatre dels quals són defectuosos. Se'n trauen tres de manera successiva i sense tornar-los a la caixa.\begin{tasks}(1)
  \task Calcula la probabilitat que cap dels tres bolígrafs tinga defectes.
  \task Calcula la probabilitat que almenys un dels tres siga defectuós.
  \task Calcula la probabilitat que només un siga defectuós.
\end{tasks}

\ex{2} \pau{Aragó, EvAU juny 2017} En una urna hi ha 2 boles blanques, 4 de negres i 5 de roges. Se'n trauen dues, una darrere l'altra, sense reemplaçament. Calcula:\begin{tasks}(1)
  \task La probabilitat que les dues siguen roges.
  \task La probabilitat que siguen de color diferent.
  \task La probabilitat que la segona bola siga roja.
  \task Si $A$ és el succés «la primera bola és roja» i $B$ el succés «les dues boles són del mateix color», són independents $A$ i $B$?
\end{tasks}

\ex{2} \pau{Andalusia, PEvAU extraordinària 2023} Una caixa conté 3 fitxes verdes, 2 de blaves i 4 de roges. Un joc consisteix a fer dues extraccions sense reemplaçament: qui trau dues fitxes blaves guanya el primer premi; qui en trau dues de verdes, el segon premi; i qui trau una blava i una d'un altre color, el tercer premi.\begin{tasks}(1)
  \task Calcula la probabilitat que un jugador aconseguisca el primer o el segon premi.
  \task Calcula la probabilitat que un jugador guanye el tercer premi.
  \task Sabent que un jugador ha obtingut premi, quina és la probabilitat que haja guanyat el tercer?
\end{tasks}

\ex{2} \pau{La Rioja, EBAU ordinària 2023} Un bombo de loteria té deu boles numerades del 0 al 9. Fem dues extraccions consecutives sense reemplaçar la primera bola. Considerem els successos $A=$«la primera bola és 0», $B=$«la primera bola és 5» i $C=$«la segona bola és més gran que la primera». Calcula $P(A)$, $P(B)$, $P(C)$, $P(A\cond C)$ i $P(B\cond C)$.

\ex{2} \pau{País Basc, EAU extraordinària 2022} Es van a sortejar 4 viatges a Nova York entre 40 persones amb una baralla de 40 cartes: es reparteix una carta a cada persona i les que reben un rei guanyen un viatge.\begin{tasks}(1)
  \task Calcula la probabilitat que guanye un viatge la primera persona que rep la carta.
  \task Calcula la probabilitat que guanye un viatge la segona persona que rep la carta.
  \task Calcula la probabilitat que cap de les dues primeres persones guanye un viatge.
\end{tasks}

\ex{2} \pau{Múrcia, EBAU juny 2017} Una urna conté tres boles numerades de l'1 al 3. Se'n trauen successivament les tres boles.\begin{tasks}(1)
  \task Calcula la probabilitat que les dues últimes boles extretes siguen senars.
  \task Determina si són independents els successos $S_1=$«ix el nombre parell abans que algun dels senars» i $S_2=$«els dos nombres senars ixen consecutivament».
\end{tasks}

\ex{2} \pau{Aragó, EvAU setembre 2019} Una empresa té 64 treballadors repartits en tres departaments. S'ha estudiat si saben anglés o no:

\begin{center}{\small\renewcommand{\arraystretch}{1.25}\begin{tabular}{l|ccc|c}
 & $\text{Admin.}$ & $\text{Producció}$ & $\text{Vendes}$ & \textbf{Total}\\\hline
$\text{Sap anglés}$ & $12$ & $30$ & $6$ & $48$\\
$\text{No en sap}$ & $4$ & $11$ & $1$ & $16$\\
\hline
\textbf{Total} & $16$ & $41$ & $7$ & $64$\\
\end{tabular}}\end{center}
\begin{tasks}(1)
  \task Triem a l'atzar un treballador. Quina és la probabilitat que sàpia anglés?
  \task Triem a l'atzar un treballador d'entre els que saben anglés. Quina és la probabilitat que siga del departament de Vendes?
  \task Siguen $A=$«el treballador és d'Administració» i $B=$«el treballador sap anglés». Són independents?
  \task Triem a l'atzar tres treballadors (sense reemplaçament). Quina és la probabilitat que siguen del mateix departament?
\end{tasks}


% =====================================================
\seccio{8}{Probabilitat total i teorema de Bayes}

\exf{1} \pau{Aragó, EvAU ordinària 2023} En una facultat d'Economia s'ofereix una assignatura en tres grups, G1, G2 i G3, que representen el 40\,\%, el 35\,\% i el 25\,\% dels estudiants, respectivament. Supera l'assignatura el 80\,\% del grup G1, el 60\,\% del G2 i el 92\,\% del G3. Calcula la probabilitat que, en triar un estudiant a l'atzar:\begin{tasks}(1)
  \task haja superat l'assignatura i siga del grup G3;
  \task no haja superat l'assignatura;
  \task haja superat l'assignatura;
  \task ni haja superat l'assignatura ni siga del grup G1;
  \task siga del grup G3, si ha superat l'assignatura.
\end{tasks}\recorda{probabilitat total: $P(S)=P(G_1)P(S\cond G_1)+P(G_2)P(S\cond G_2)+P(G_3)P(S\cond G_3)$. Bayes: $P(G_i\cond S)=\dfrac{P(G_i)P(S\cond G_i)}{P(S)}$.}

\ex{1} \pau{Andalusia, PEvAU setembre 2018} Un hotel té tres rentadores industrials, $L_1$, $L_2$ i $L_3$. $L_1$ fa el 50\,\% de les rentades, $L_2$ el 30\,\% i $L_3$ la resta. $L_1$ produeix un 5\,\% de rentades defectuoses, $L_2$ un 15\,\% i $L_3$ un 20\,\%. Es tria a l'atzar una rentada.\begin{tasks}(1)
  \task Calcula la probabilitat que no siga defectuosa.
  \task Calcula la probabilitat que l'haja feta $L_1$, sabent que ha sigut defectuosa.
\end{tasks}

\ex{1} \pau{Galícia, ABAU setembre 2018} En una empresa, el 30\,\% dels empleats són dones i el 70\,\% restant, homes. De les dones, el 80\,\% té contracte indefinit, mentre que dels homes només el 70\,\% en té.\begin{tasks}(1)
  \task Calcula el percentatge de persones de l'empresa que tenen contracte indefinit.
  \task Si un empleat té contracte indefinit, quina és la probabilitat que siga dona?
  \task Són independents els successos «ser home» i «tenir contracte indefinit»?
\end{tasks}

\ex{1} \pau{Galícia, ABAU juny 2018} El 30\,\% dels estudiants de secundària d'un institut juga a bàsquet. Entre els que juguen a bàsquet, el 40\,\% juga també a tenis; entre els que no juguen a bàsquet, una quarta part juga a tenis. Triem un estudiant a l'atzar.\begin{tasks}(1)
  \task Quina és la probabilitat que practique tots dos esports?
  \task Quina és la probabilitat que jugue a tenis?
  \task Són independents els successos «jugar a tenis» i «jugar a bàsquet»?
\end{tasks}

\ex{1} \pau{Andalusia, PEvAU extraordinària 2020} Una empresa fabrica bombetes LED i halògenes. El 5\,\% de les LED i el 2\,\% de les halògenes ixen defectuoses. Es tria a l'atzar una bombeta d'una caixa que conté 40 bombetes LED i 10 d'halògenes.\begin{tasks}(1)
  \task Calcula la probabilitat que la bombeta triada no siga defectuosa.
  \task Calcula la probabilitat que la bombeta siga LED, sabent que és defectuosa.
\end{tasks}

\ex{2} \pau{Andalusia, PEvAU ordinària 2023} En una base de dades de correus electrònics s'ha observat que el 20\,\% dels correus rebuts són brossa (\emph{spam}). A més, la paraula «lottery» apareix en el 40\,\% dels correus brossa i en el 0,6\,\% dels que no ho són.\begin{tasks}(1)
  \task Calcula la probabilitat que un correu triat a l'atzar en què apareix la paraula «lottery» siga brossa.
  \task Calcula la probabilitat que un correu triat a l'atzar en què no apareix la paraula «lottery» no siga brossa.
  \task Si un correu s'etiqueta com a brossa quan apareix la paraula «lottery» i com a no brossa quan no apareix, calcula la probabilitat que un correu s'etiquete incorrectament.
\end{tasks}

\ex{2} \pau{Navarra, EvAU extraordinària 2022} Es considera que el 4\,\% dels esportistes federats consumeix alguna substància no permesa. Una nova prova de detecció identifica correctament el consum en el 92\,\% dels casos, però si un esportista no en consumeix, la prova dona positiu en el 10\,\% dels casos. Es tria un esportista a l'atzar. Calcula (amb 4 decimals):\begin{tasks}(1)
  \task La probabilitat que obtinga un resultat positiu.
  \task La probabilitat que siga consumidor i la prova done negatiu.
  \task La probabilitat que no siga consumidor, sabent que la prova ha donat negatiu.
\end{tasks}

\ex{2} \pau{Galícia, ABAU extraordinària 2023} En una urna A hi ha 8 boles verdes i 6 de roges, i en una urna B, 4 de verdes i 5 de roges. Es llança un dau: si ix un nombre menor que 3 es trau una bola de l'urna A, i si no, de l'urna B.\begin{tasks}(1)
  \task Calcula la probabilitat que la bola extreta siga roja.
  \task Sabent que s'ha extret una bola verda, quina és la probabilitat que siga de l'urna A?
  \task Són independents els successos «traure bola roja» i «la bola és de l'urna A»?
\end{tasks}

\ex{2} \pau{Canàries, EBAU extraordinària 2022} Tenim tres urnes iguals: la primera té 6 boles blanques i 4 de negres, la segona 5 de blanques i 2 de negres, i la tercera 4 de blanques i 7 de negres.\begin{tasks}(1)
  \task Es trau una bola d'una urna triada a l'atzar. Fes un diagrama amb les probabilitats dels resultats possibles.
  \task Calcula la probabilitat de traure una bola negra.
  \task S'ha extret una bola i és blanca. Calcula la probabilitat que siga de la primera urna.
\end{tasks}

\ex{2} \pau{Múrcia, EBAU extraordinària 2021} Hi ha tres caixes: la primera conté 4 boles blaves i 6 de blanques; la segona, una bola blava i 5 de blanques; la tercera, 3 de blaves i 5 de blanques. Agafem una bola a l'atzar d'una de les caixes, també triada a l'atzar.\begin{tasks}(1)
  \task Calcula la probabilitat que la bola siga blava.
  \task Si la bola és blanca, calcula la probabilitat que siga de la primera caixa.
\end{tasks}

\ex{1} \pau{Navarra, EvAU setembre 2018} Dos estudiants construeixen un dau amb 5 cares roges i una de blava, posen 7 boles verdes i 3 de negres en una caixa A, i 6 de negres i 2 de verdes en una caixa B. Llancen el dau: si ix roig, traguen una bola de la caixa A, i si ix blau, de la B.\begin{tasks}(1)
  \task Calcula la probabilitat que la bola siga verda.
  \task Calcula la probabilitat que la cara del dau siga blava, sabent que la bola no és verda.
\end{tasks}

\ex{1} \pau{Aragó, EvAU extraordinària 2024} A Espanya, el 30\,\% de la població té menys de 30 anys, el 50\,\% entre 30 i 65 anys, i el 20\,\% més de 65. Un estudi afirma que té telèfon mòbil el 70\,\% dels menors de 30 anys, el 95\,\% dels que tenen entre 30 i 65, i el 50\,\% dels majors de 65.\begin{tasks}(1)
  \task Es tria una persona a l'atzar. Calcula la probabilitat que tinga més de 65 anys i tinga mòbil.
  \task Quina és la probabilitat que una persona triada a l'atzar tinga mòbil?
  \task Triem una persona a l'atzar i té mòbil. Quina és la probabilitat que tinga més de 65 anys?
  \task Triem una persona de cada grup d'edat. Quina és la probabilitat que totes tres tinguen mòbil? (Es pot suposar independència.)
\end{tasks}

\ex{1} \pau{Galícia, ABAU setembre 2018} En una població, de cada 200 consumidors d'una beguda isotònica, 60 consumeixen la marca A, 50 la marca B i la resta, la C. A més, són joves el 30\,\% dels consumidors d'A, el 20\,\% dels de B i el 40\,\% dels de C.\begin{tasks}(1)
  \task Si se selecciona a l'atzar un consumidor, quina és la probabilitat que siga jove?
  \task Si s'ha seleccionat un jove, troba la probabilitat que consumisca la marca B.
  \task Són independents els successos «ser jove» i «consumir la marca A»?
\end{tasks}

\ex{2} \pau{Andalusia, PEvAU extraordinària 2022} En una regió hi ha tres universitats, A, B i C. Dels estudiants que van acabar l'any passat, el 60\,\% procedien de la A, el 30\,\% de la B i la resta de la C. La probabilitat que un estudiant de la A no trobe treball a la regió és 0,4, i per a un de la B, 0,5.\begin{tasks}(1)
  \task Si la probabilitat que un estudiant no trobe treball a la regió és 0,395, determina la probabilitat que un estudiant de la universitat C sí que en trobe.
  \task Calcula la probabilitat que un estudiant que no ha trobat treball a la regió procedisca de la universitat A o de la B.
\end{tasks}

\ex{2} \pau{Andalusia, PEvAU extraordinària 2020} Una urna conté 6 boles roges i 4 de blaves. Es trau una bola a l'atzar i es reemplaça per sis boles de l'altre color. A continuació, es trau una segona bola.\begin{tasks}(1)
  \task Calcula la probabilitat que la segona bola siga roja.
  \task Si sabem que la segona bola és blava, quina és la probabilitat que la primera també ho haja sigut?
\end{tasks}

\ex{2} \pau{País Basc, EAU ordinària 2024} L'Asier té una urna amb 4 boles verdes i 2 de roges. Llança una moneda: si ix cara, trau una bola de l'urna, i si ix creu, en trau dues sense reemplaçament.\begin{tasks}(1)
  \task Quina és la probabilitat que l'Asier haja tret dues boles roges?
  \task Calcula la probabilitat que no haja tret cap bola roja.
  \task Calcula la probabilitat que haja tret almenys una bola verda.
  \task Calcula la probabilitat que haja eixit cara, sabent que almenys una bola és verda.
\end{tasks}

\ex{2} \pau{País Basc, EAU extraordinària 2024} En una caixa hi ha una bola roja i una de blava. Es trau una bola i, abans de traure'n una segona, es torna a la caixa la primera bola extreta i se n'afigen dues més del mateix color. A continuació es trau una segona bola.\begin{tasks}(1)
  \task Calcula la probabilitat que la segona bola siga roja si la primera ha sigut blava.
  \task Calcula la probabilitat que la segona bola siga blava.
  \task Si la segona bola ha sigut blava, quina és la probabilitat que la primera haja sigut roja?
\end{tasks}

\ex{2} \pau{Madrid, EvAU ordinària 2022} D'una baralla de pòquer de 52 cartes (13 de cada pal: diamants, cors, piques i trèvols) s'elimina una carta a l'atzar, sense mirar-la. Després es tria a l'atzar una carta de les que queden i s'observa.\begin{tasks}(1)
  \task Calcula la probabilitat que la carta observada siga de diamants.
  \task Si la carta observada no és de diamants, calcula la probabilitat que la carta eliminada tampoc no ho fóra.
\end{tasks}

\vfill
\begin{center}\color{muted}\small\faCheckCircle[regular]\ \ Ànim! Comprova els teus resultats amb les eines interactives dels apunts de la unitat.\end{center}

\end{document}
